The Moon and planets are shown in the projected sky, if they are above the
horizon. The Milky Way and other diffuse objects are not shown. In the
planetarium's projection, azimuth starts at the North Cardinal Point, which
is assigned 0 degrees; the East Cardinal Point has an azimuth of 90 degrees,
and so on.

\begin{description}
\item[Question 1 (4 points)] Write the names of the zodiacal constellations
  (in either Latin or the standard IAU abbreviation) that are partially or
  completely above the horizon.
\item[Question 2 (4 points)] There is a bright star (brighter than 1.3 mag)
  located at each of the Alt-Az coordinates given below. Give their Bayer
  designation.
  \begin{description}
  \item[(a)] Star 1: Azimuth: $254^\circ 23'$; Altitude: $68^\circ 22'$
  \item[(b)] Star 2: Azimuth: $113^\circ 24'$; Altitude: $24^\circ 30'$
  \end{description}
\item[Question 3 (2 points)]~
  \begin{description}
  \item[(a)] The alpha star of a boreal constellation (i.e. northern
    hemisphere constellation) lies almost on the local meridian. Give the
    Bayer designation of this star.
  \item[(b)] Give the Messier Catalogue number, or New General Catalogue
    designation, of a nebula that is located almost at the intersection of
    the Ecliptic and the local meridian.
  \end{description}
\end{description}
